\section{Introduction}
\label{sec:introduction}

\subsection{Background}

Housing is the largest asset most households own and a central channel
through which monetary policy, credit conditions and financial stability
interact. Episodes in which house prices detach from fundamentals have
preceded some of the most costly downturns of recent decades, and they have
rarely been confined to one country: booms and busts in housing markets tend
to be synchronised internationally. Monitoring them therefore calls for
evidence that is comparable across countries, grounded in a consistent
methodology, and timely enough to matter while developments are still
unfolding.

The International House Price Database and the International Housing
Observatory were built to provide that evidence. The database supplies a
long, comparable cross-country panel of house prices and incomes; the
Observatory applies a common set of econometric tools to it each quarter and
publishes the results openly. Together, they aim to be a reference point for
researchers, policymakers and the public following developments in
international housing markets.

\subsection{The International House Price Database}

Comparing house prices across countries is harder than it first appears.
National statistical agencies, central banks and private data providers each
construct their indices for their own purposes, and those purposes differ:
indices vary in geographic coverage, in the vintage and type of dwelling
priced, in how they control for quality change, in reporting frequency, in
base year, and in whether and how they are seasonally adjusted. Differences
of this kind are large enough that a naive cross-country panel assembled
from national sources will confound genuine price dynamics with measurement
convention \citep{mack2011}.

The Dallas Fed International House Price Database (IHPD) exists to address
exactly this problem. For each economy it selects the national series most
consistent with a common benchmark, extends that series back to 1975\,Q1
using historical sources, and then applies a uniform set of transformations
--- frequency conversion, splicing, seasonal adjustment, rebasing --- so that
the resulting panel is comparable across countries. It pairs each house price
series with a matching measure of personal disposable income, so that prices
need not be read in isolation from the capacity of households to pay them.

\subsection{The International Housing Observatory}

The International Housing Observatory is a collaborative research initiative
of the Federal Reserve Bank of Dallas (Globalization Institute), Lancaster
University Management School (Economics Department), and the University of
Alicante (Department of Fundamentals of Economic Analysis). Its public
platform, at \url{https://int.housing-observatory.com/}, is intended as a
place to follow developments in international housing markets as they
happen: each new data release is passed through a fixed econometric pipeline
and published as an updated set of exuberance indicators, charts and reports.

In practice the Observatory offers four things: an interactive dashboard
showing, economy by economy, whether and when prices have behaved
explosively; a quarterly monitoring report that places those results in
economic context; a nowcast of real house prices ahead of the official
release; and, through the companion Housing Fever site, an assessment of
whether price movements are accounted for by fundamentals. These outputs are
described in Section~\ref{sec:website}.

The Observatory's role is therefore complementary to that of the database.
The Dallas Fed constructs and publishes the data; the Observatory maintains
the statistical machinery applied to it, the open-source software that
implements that machinery, and the public-facing outputs through which the
results are disseminated.

\subsection{Purpose of This Report}

The Observatory's data, methods and outputs are currently documented across
several places: the working paper describing the database, a set of published
econometric papers, two \textsf{R} package manuals, and the methodology pages
of the website. There is no single document that describes the system as a
whole, and no citable reference for the platform itself. This report is
intended to be that document.

It is deliberately \emph{not} a restatement of
\citet{mack2011}. That working paper remains the authoritative description of
how the IHPD is built, is maintained by the team that builds it, and is
revised whenever the construction methodology changes. Section~\ref{sec:database}
summarises its methodology at the level of detail needed to interpret the
data correctly and to understand the limitations that follow from it, and
refers the reader there for the remainder --- in particular for the
country-by-country source documentation, which this report does not attempt
to duplicate.

\subsection{Structure of the Report}

The remainder of the report is organised as follows:

\begin{enumerate}
  \item \textbf{The database} (Section~\ref{sec:database}) --- coverage,
        variable definitions, and the construction methodology of
        \citet{mack2011}, including the consequences of that methodology for
        how the early part of the sample should be read.

  \item \textbf{Data access} (Section~\ref{sec:data-access}) --- the structure
        of the quarterly release files, the Observatory's own JSON data API, and the quarterly update workflow.

  \item \textbf{Econometric methodology} (Section~\ref{sec:methodology}) ---
        the GSADF test and BSADF date-stamping procedure
        \citep{phillips2015a, phillips2015b}, the PSY-IVX decomposition
        \citep{shi2021}, and the mixed-frequency dynamic factor model used to
        nowcast real house prices ahead of the official release.

  \item \textbf{Platform and outputs} (Section~\ref{sec:website}) --- the
        interactive dashboard, the quarterly monitoring reports, the nowcast
        reports, the Housing Fever companion platform, and the supporting
        open-source software.

  \item \textbf{Planned extensions} (Section~\ref{sec:roadmap}) --- additions
        to the database and to the econometric toolkit that are agreed or
        under consideration.
\end{enumerate}

\noindent Country coverage and the sources of country-level documentation are
given in Appendix~\ref{app:countries}. The version history of this report is
in Appendix~\ref{app:changelog}.
